% Cover Letter — IEEE Transactions on Image Processing submission
% NOTE: TIP states that submissions lacking a detailed cover letter
%       may be returned without review.
% Compile: pdflatex mic-cover-letter-ieee-tip.tex
\documentclass[12pt]{article}

\usepackage[margin=1in]{geometry}
\usepackage{xcolor}
\usepackage{hyperref}
\usepackage{parskip}
\usepackage{microtype}

\hypersetup{
  colorlinks=true,
  urlcolor=blue!70!black,
}

\pagestyle{empty}

\begin{document}

\begin{flushright}
Kuldeep Singh\\
Innovaccer Inc.\\
\href{mailto:kuldeep.singh@innovaccer.com}{kuldeep.singh@innovaccer.com}\\
ORCID: 0009-0004-8476-0118\\[6pt]
May 2, 2026
\end{flushright}

\bigskip

Editor-in-Chief\\
\textit{IEEE Transactions on Image Processing}

\bigskip

\noindent Dear Editor-in-Chief,

\medskip

I am submitting the manuscript \textbf{``A 16-Bit-Native Entropy Coding
Pipeline for Lossless Compression of High-Bit-Depth Images''} for consideration
as a Full Paper in \textit{IEEE Transactions on Image Processing}.

\medskip
\noindent\textbf{Summary and core contribution.}
This paper addresses a fundamental mismatch in lossless image compression:
most entropy coders are optimized for 8-bit byte streams, yet high-bit-depth
images (medical scanners, scientific instruments, HDR cameras) store pixel
data at 10--16 bits per sample. Applying a byte-oriented coder to 16-bit data
requires splitting each pixel into a byte pair before entropy coding, losing
the mutual information between the two bytes. I quantify this structural cost
as 0.3--1.1 bits per symbol after spatial prediction, and present MIC---a
lossless codec that eliminates it by coding 16-bit residuals as atomic symbols.

MIC combines spatial prediction, a 16-bit-native run-length representation,
and an extended Finite State Entropy (FSE/ANS) coder supporting up to
65,535-symbol alphabets. This contrasts with Zstandard's FSE capped at 4,096
symbols and Huffman coding, which becomes impractical above roughly 4,096
symbols. The codec also includes a four-state interleaved decoder that exploits
instruction-level parallelism (ILP), yielding a geometric-mean 68\% isolated
decompression speedup on ARM64 (43\% on AMD64) with no compressed-size
penalty.

\medskip
\noindent\textbf{Relevance to IEEE TIP.}
The manuscript makes contributions in image coding, entropy coding, and
decoder algorithm design---all within the core scope of IEEE TIP. The work:
\begin{itemize}
  \item proposes and analyzes a new entropy coding approach for large symbol
  alphabets arising in high-bit-depth image residual streams;
  \item provides an information-theoretic analysis (mutual information of
  byte-split residuals) that explains and predicts the compression advantage;
  \item evaluates a multi-state decoder design that improves throughput via
  instruction-level parallelism, a technique of broad relevance to image codec
  design.
\end{itemize}
The evaluation uses 21 de-identified high-bit-depth images (nine modalities,
sizes up to $4774 \times 3064$ pixels), with fair in-process comparisons
against High-Throughput JPEG~2000 (HTJ2K) and JPEG-LS, two codecs within
IEEE TIP's reference space.

\medskip
\noindent\textbf{Suggested EDICS categories.}
\texttt{COD-CMP} (Compression); \texttt{COD-ENT} (Entropy Coding);
\texttt{ANA-BIO} (Biomedical image analysis) as a secondary category.

\medskip
\noindent\textbf{Originality.}
This work has not been submitted to, or published in, any other journal.
The manuscript is an original work; no prior conference or workshop version
has been published.

\medskip
\noindent\textbf{Open-source reproducibility.}
The complete implementation, benchmark harness, and test data references
are at \url{https://github.com/pappuks/medical-image-codec}. All reported
results are reproducible using the commands in the repository.

\medskip
\noindent\textbf{Ethics.}
All 21 test images are from publicly available, de-identified DICOM datasets
(NEMA WG-04 compsamples, NEMA 1997~CD, Clunie Breast Tomosynthesis Case~1).
No human subjects research was conducted; no IRB approval is required.

\medskip
\noindent\textbf{Conflict of interest.}
The author declares no competing financial interests or personal relationships
that could have influenced this work.

\medskip
\noindent\textbf{Suggested reviewers.}
Researchers active in ANS/FSE entropy coding, image coding standards (HTJ2K,
JPEG-LS), or high-throughput codec design would be well placed to review
this work. Authors of OpenJPH and CharLS are compared against in the paper
and may have a conflict of interest at the Editor's discretion.

\bigskip

I confirm this is a single-author submission and that the manuscript complies
with IEEE Transactions on Image Processing formatting requirements, including
the 13-page initial-submission limit and the 150--250 word abstract.

\bigskip

\noindent Sincerely,

\bigskip\bigskip

\noindent Kuldeep Singh\\
Innovaccer Inc.\\
\href{mailto:kuldeep.singh@innovaccer.com}{kuldeep.singh@innovaccer.com}\\
ORCID: 0009-0004-8476-0118\\
\url{https://github.com/pappuks/medical-image-codec}

\end{document}
